\chapter{概念界定}


\section{习题课}

习题课是教师按照课标中所描述的学生需达到的能力要求以及教材的内容，在课堂上主要开展习题讲解、题型训练、知识归纳等教学活动的典型课型，其中习题讲解是引导学生突破学习中重难点的重要手段，题型训练是巩固知识和检查问题的主要环节，知识归纳是使知识系统化、模块化的主要措施。不同于新授课，新授课强调知识的传授与接收，注重发挥教师的主导作用，习题课强调以学生已学过的知识为基础，利用习题进行知识的巩固、检验、应用，注重发挥学生的主体作用，渗透思想方法与实践精神。



\section{数学探究性学习}

数学探究性学习通常从课程论和教与学两个角度叙述，基于本文的探讨主题，笔者主要从教与学这一角度进行叙述。

范文贵在《数学探究学习内涵与策略研究》中提到：“数学探究性学习是在教师的组织指导下，基于数学问题解决所需的情境和途径，学生充分利用各种数学探究课程资源，通过实物操作实验、数学思想实验、数据统计分析、信息技术实验等，收集与研究问题相关的各种信息，学生提出研究问题、形成猜想、利用特例进行实验验证、进行演绎证明、反思(推广拓展) \cite{范文贵2005数学探究学习内涵与策略研究}。”

宁连华在《数学探究学习探究》中阐述道：“概括地说，数学探究学习是指学生自己或合作共同体针对要学习的概念、原理、法则或要解决的数学问题主动地思考、探索的学习活动，强调的是一种主动参与的学习方式\cite{宁连华2004数学探究学习研究}。”他强调数学探究性学习拥有一般的探究学习的特征，主要形式是解题活动，比一般的探究学习更有思维性。

基于以上学者研究数学探究性学习所得出的定义与内涵，本文将数学探究性学习定义为：在教师的指导下，学生在数学领域的问题情境中，以自主学习和合作交流为途径，通过主动参与探究的过程并得出相应的结论，获取到问题的解决方案。



\section{数学习题课中的探究性学习}

众多数学教育家都对在数学学习中需要有大量的数学习题训练给予了极高的重视，大多学者强调的并非简单机械地盲目模仿训练，而是给高质量的、思考性强的问题留时间探索。宁连华提出：找出适合探究的习题至关重要，这类习题介于一般习题与数学研究之间，也是数学学习中一类及其重要的问题。但与常规的数学习题相比，此类问题并不存在本质区别，常规习题可以引导延伸到更加本质、深入的可待研究的层面上\cite{宁连华2004数学学科探究学习的特征及其指导策略}。

基于以上学者的观点与本文研究内容，本文将数学习题课中的探究性学习定义为：在数学习题课的课堂上，在教师的引导及指导下，学生以数学学习中的常规习题为基础探究内容，在习题所反应的可待探究的层面主动参与提出问题、自主学习、合作交流等探究活动，获取到习题的解决方案，得到习题延伸的启示。